% HAND-WRITTEN fragment -- NOT generated by maestro2tex. Input from
% AnalogChapter.tex after note_eis_excitation.tex, last in the impedance
% section.  ADDED 2026-08-26 (owner request): "a few equations for how firmware
% turns the sequence of digitised values into R and C, just the equations".
%
% Nothing here is a new measurement.  Every equation is the algebra already
% implied by statements elsewhere in this chapter, written out:
%   (1) the demodulator: note_eis_intro, "an I/Q multiply and accumulate over an
%       integer number of excitation periods, no window, no detrending".  The
%       normalisation 2/MN is what makes X-hat the AMPLITUDE phasor: for
%       c_n = A cos(2 pi n/N + phi), sum c_n exp(-j 2 pi n/N) = (MN/2) A e^{j phi}.
%   (2) the transimpedance sign: note_ampab_tia sss:ampab-tia states
%       V_OUT = v_cm - I_WE R_f for A2, so in the fundamental bin (where the dc
%       common mode does not appear) I_WE-hat = -V_OUT-hat / R_f.  The leading
%       minus in eq:eis-z-rf is that inversion and nothing else.
%   (3) the LSB cancels because Z is a ratio of two phasors taken from the same
%       converter; likewise the 4/pi fundamental of the square (note_eis_excitation,
%       "the demodulator is ratiometric and the excitation's shape cancels").
%   (4) the ratiometric form is the architecture this section requires four times
%       over (feedback pole, code-pair error, quantisation bias, sigma = 11 % on
%       R_f).  It is architecture "C" of the source campaign.
%   (5) the cell model is note_models: R_ct || C_dl in series with R_u.
%       Z - R_u = R_ct/(1 + j w R_ct C_dl), so 1/(Z - R_u) = 1/R_ct + j w C_dl.
%       That identity is the whole reason the extraction is two divisions and no
%       fit: R_ct falls out of the real part and C_dl out of the imaginary part
%       of ONE inverted mid-band point.
%
% POTENTIAL CHANNEL, deliberate wording.  V_P-hat is written as "whatever the
% converter multiplexer presents as the potential channel" rather than as the RE
% tap, because on the shipping cell it CANNOT be the RE tap: ADC_MUX has no
% reference-electrode input (sss:pixtop-adc), so slot 2, the commanded potential,
% is what firmware actually demodulates.  That is architecture "D" of the source
% campaign, and the gap to "A" is the 0.02-0.2 % of magnitude quoted in
% note_eis_accuracy on this axis and the up-to-3.5 % of note_eis_excitation on
% the excitation axis.  Do NOT tighten this to "the reference electrode": it
% would state a connection the die does not have.
%
% No number is introduced here.  Every figure quoted is a back-reference.
% Rendered prose carries no em or en dashes.

\subsubsection{From Codes to Resistance and Capacitance}\label{sss:eis-firmware}

Firmware holds two code sequences per frequency, the potential channel and the
transimpedance output, taken on a grid of \(N\) samples per excitation period
over \(M\) whole periods. One multiply and accumulate per channel reduces each
to its fundamental phasor,
\begin{equation}\label{eq:eis-demod}
  \hat{X} \;=\; \frac{2}{MN}\sum_{n=0}^{MN-1} c_n\,
                \mathrm{e}^{-\mathrm{j}2\pi n/N},
  \qquad \mathrm{j}^2 = -1,
\end{equation}
giving \(\hat{V}_{\mathrm{P}}\) and \(\hat{V}_{\mathrm{OUT}}\). The converter LSB
is never applied: it cancels in every ratio below, as does the \(4/\pi\)
fundamental of the two-level excitation.

Since A2 holds the working electrode at the common mode across \(R_f\)
(Section~\ref{sss:ampab-tia}), the electrode current in that bin is
\(\hat{I}_{\mathrm{WE}} = -\hat{V}_{\mathrm{OUT}}/R_f\), so with a per-chip dc
calibration of the feedback resistance
\begin{equation}\label{eq:eis-z-rf}
  \hat{Z} \;=\; \frac{\hat{V}_{\mathrm{P}}}{\hat{I}_{\mathrm{WE}}}
          \;=\; -R_f\,\frac{\hat{V}_{\mathrm{P}}}{\hat{V}_{\mathrm{OUT}}},
\end{equation}
the minus sign being the transimpedance inversion. Measuring a known
\(Z_{\mathrm{cal}}\) at the same frequency and gain code removes \(R_f\)
altogether, and with it the feedback pole and the quantisation bias:
\begin{equation}\label{eq:eis-z-ratio}
  \hat{Z} \;=\; Z_{\mathrm{cal}}\,
    \frac{\hat{V}_{\mathrm{P}}\,\hat{V}^{\,\mathrm{cal}}_{\mathrm{OUT}}}
         {\hat{V}_{\mathrm{OUT}}\,\hat{V}^{\,\mathrm{cal}}_{\mathrm{P}}}.
\end{equation}

Against the cell model of Section~\ref{ss:analog-models},
\begin{equation}\label{eq:eis-cell}
  Z(\omega) \;=\; R_u \;+\; \frac{R_{\mathrm{ct}}}
                                 {1 + \mathrm{j}\omega R_{\mathrm{ct}} C_{\mathrm{dl}}},
\end{equation}
the extraction is two divisions and no curve fit. Take \(R_u\) as
\(\Re\,\hat{Z}\) at the top of the band, where the parallel branch has
collapsed, then invert one mid-band point:
\begin{equation}\label{eq:eis-extract}
  Y \;=\; \frac{1}{\hat{Z} - R_u}
    \;=\; \frac{1}{R_{\mathrm{ct}}} + \mathrm{j}\omega C_{\mathrm{dl}},
  \qquad
  R_{\mathrm{ct}} = \frac{1}{\Re\,Y},
  \qquad
  C_{\mathrm{dl}} = \frac{\Im\,Y}{\omega}.
\end{equation}
Both readings come from the same inversion, so a single frequency inside
\SI{0.1}{\hertz} to \SI{5.3}{\kilo\hertz} yields both, and the accuracy of each
is the accuracy of \(\hat{Z}\) at that frequency, which the calibration paragraphs above bound.

\noindent The potential phasor \(\hat{V}_{\mathrm{P}}\) is whatever the
converter multiplexer presents as that channel. On this cell that is the commanded potential,
not a reference-electrode tap, which no multiplexer slot carries
(the converter readout note records the slot map); the difference is worth
\SIrange{0.02}{0.2}{\percent} of magnitude here and up to \SI{3.5}{\percent} on
the excitation (the excitation and averaging note carries that figure).
